% ch6_b 第6章 §6.3–§6.4 (书页 263–274 / p278–p289)
% 块界说明：书页 262 以完整段落结束（"…we have solved the compact U(1) gauge theory on a torus."），
% 本块起始段为新段，无需 %JOIN%。本块为第 6 章章末块，至书页 274（p289）自然收束，无 TAIL。

下面我们转向 XY 模型及其本征态。我们展开
\begin{equation*}
\theta(t,\pmb{x})=\theta_{0}(t)+\frac{2\pi}{L_{i}}m_{i}x^{i}+\sum_{\pmb{k}}\lambda_{\pmb{k}}(t)\,\mathrm{e}^{\mathrm{i}\pmb{k}\cdot\pmb{x}}.
\end{equation*}
这里 $m_{1}$ 与 $m_{2}$ 分别是 $\theta$ 在 $x^{1}$ 与 $x^{2}$ 方向上的缠绕数。XY 模型的拉格朗日量现在具有如下形式：
\begin{equation*}
L=\frac{\chi}{2}\left(L_{1}L_{2}\dot{\theta}_{0}^{2}-\frac{(2\pi)^{2}L_{2}}{L_{1}}m_{1}^{2}-\frac{(2\pi)^{2}L_{1}}{L_{2}}m_{2}^{2}+L_{1}L_{2}\sum_{\pmb{k}}\left(|\dot{\lambda}_{\pmb{k}}|^{2}-\pmb{k}^{2}|\lambda_{\pmb{k}}|^{2}\right)\right)
\end{equation*}
我们看到，XY 模型由一群振子（由 $\lambda_{\pmb{k}}$ 描述）和一个圆上的自由粒子（由 $\theta_{0}$ 描述）构成。XY 模型的本征态由 $(N,m_{i},n_{\pmb{k}})$ 标记，其能量为
\begin{equation*}
E=\frac{1}{2\chi L_{1}L_{2}}N^{2}+\frac{\chi(2\pi)^{2}L_{2}}{2L_{1}}m_{1}^{2}+\frac{\chi(2\pi)^{2}L_{1}}{2L_{2}}m_{2}^{2}+\sum_{\pmb{k}}|\pmb{k}|n_{\pmb{k}} \tag{6.3.8}
\end{equation*}
我们看到，如果
\begin{equation*}
\chi=\left(\frac{qg}{2\pi}\right)^{2}
\end{equation*}
并且在标记 U(1) 理论与 XY 模型的本征态时把 $n_{i}$ 认作 $\epsilon_{ij}m_{j}$，那么 XY 模型与 U(1) 规范理论就具有完全相同的谱。本征态还携带确定的总动量 $\pmb{K}$：对 XY 模型的态 $|N,m_{i},n_{\pmb{k}}\rangle$ 或 U(1) 理论的态 $|N,n_{i},n_{\pmb{k}}\rangle$，
\begin{equation*}
K_{i}=Nm_{i}\frac{2\pi}{L_{i}}+\sum_{\pmb{k}}k_{i}n_{\pmb{k}}=N\epsilon_{ij}n_{j}\frac{2\pi}{L_{i}}+\sum_{\pmb{k}}k_{i}n_{\pmb{k}}.
\end{equation*}
从 XY 模型这一侧我们看到，标记 $N$ 具有物理意义：它是玻色子的总数减去平衡时的玻色子数，即 $N=N_{tot}-N_{0}$。因此，$\frac{1}{\chi}\dot{\theta}$ 与 $\frac{q}{2\pi}b$ 可以看作玻色子数密度减去平衡密度：
\begin{equation*}
\rho-\rho_{0}=\frac{q}{2\pi}b=\frac{1}{\chi}\dot{\theta}
\end{equation*}
U(1) 理论中的约束 $\partial_{0}b-\epsilon_{ij}\partial_{i}e_{j}=0$ 可以看作守恒流的连续性方程，即 $\partial_{0}\rho+\partial_{i}j_{i}=0$。于是，若把 $\frac{q}{2\pi}b$ 认作玻色子数密度，则 U(1) 理论中的玻色子数流密度由下式给出：
\begin{equation*}
j_{i}=-\frac{q}{2\pi}\epsilon_{ij}e_{j}
\end{equation*}

XY 模型的运动方程，即 $\partial_{0}^{2}\theta-\partial_{x}^{2}\theta=0$，同样可以看作守恒流的连续性方程。我们发现，在 XY 模型中同一个流由下式给出：
\begin{equation*}
j_{i}=-\frac{1}{\chi}\partial_{i}\theta
\end{equation*}
我们发现，U(1) 理论中的场与 XY 模型中的场之间存在一个简单关系：
\begin{equation*}
\frac{q}{2\pi}b=\frac{1}{\chi}\dot{\theta}, \qquad\qquad \frac{q}{2\pi}\epsilon_{ij}e_{j}=\frac{1}{\chi}\partial_{i}\theta. \tag{6.3.9}
\end{equation*}
事实上，把式 (6.3.9) 代入式 (6.3.2)，就会把 U(1) 规范拉格朗日量变成 XY 模型的拉格朗日量 (6.3.1)。这是看出 2+1 维中 U(1) 规范理论与 XY 模型之间对偶关系最直接的途径。

U(1) 规范理论可以与电荷耦合。设 $J_{\mu}$ 为 U(1) 电荷的密度与流密度，则耦合后的拉格朗日量具有如下形式：
\begin{equation*}
\mathcal{L}=\frac{1}{2g^{2}}(e^{2}-b^{2})-J_{\mu}a_{\mu}
\end{equation*}
对于位于 $\pmb{x}=0$ 处的点电荷，我们有 $J_{0}=q\delta(\pmb{x})$，并且从运动方程得到
\begin{equation*}
\frac{1}{g^{2}}\pmb{\partial}\cdot\pmb{e}=q\delta(\pmb{x})
\end{equation*}
即
\begin{equation*}
\pmb{e}=\frac{g^{2}q\pmb{x}}{2\pi\pmb{x}^{2}}
\end{equation*}
上面的 $\pmb{e}$ 对应于 XY 模型中的一个环形流：
\begin{equation*}
\partial_{i}\theta=\frac{\chi q}{2\pi}\frac{g^{2}q}{2\pi\pmb{x}^{2}}\epsilon_{ij}x^{j}=\frac{\epsilon_{ij}x^{j}}{\pmb{x}^{2}}
\end{equation*}
我们看到，U(1) 规范理论中的量子化电荷正是 XY 模型中的量子化涡旋。

\subsection*{问题 6.3.1}

证明式 (6.3.6)。（提示：你会发现展开式 (6.3.4) 很有用。）

\subsection*{问题 6.3.2}

将 1+1 维的紧致 U(1) 规范理论量子化，假设空间是长度为 $L$ 的一个圆环，电荷量子为 $q$。求标记所有能量本征态的整数集合，并求出这些本征态的能量。

\subsection{1+2 维紧致 U(1) 规范理论的禁闭}

\begin{keypoints}
\item 1+2 维 U(1) 规范理论中的瞬子效应赋予规范玻色子有限的能隙，并导致 U(1) 电荷之间的禁闭。
\item 瞬子效应可以用对偶 XY 模型中的一个 $\cos(\theta)$ 项来描述。
\end{keypoints}

这里我们考虑虚时间中的 U(1) 规范理论：
\begin{equation*}
\mathcal{L}_{U(1)}=\frac{1}{2g^{2}}\left(e^{2}+b^{2}\right) \tag{6.3.10}
\end{equation*}
形式上看，U(1) 规范理论像是一个具有无能隙激发的自由理论。然而事情并非如此简单。对于具有有限截断标度的紧致 U(1) 规范理论，理论中含有瞬子。瞬子效应将使 U(1) 规范理论成为一个相互作用理论。我们将会看到，这种相互作用会以剧烈的方式影响 U(1) 规范理论的低能性质。

什么是瞬子？位于 $x^{\mu}=0$ 处的瞬子由如下位形描述：
\begin{equation*}
b(x^{\mu})=\frac{1}{2q}\frac{x^{0}}{|\pmb{x}|^{3}}, \qquad e_{i}(x^{\mu})=\frac{1}{2q}\frac{x^{i}}{|\pmb{x}|^{3}}
\end{equation*}
上述瞬子的设计意图是使通量改变 $2\pi/q$，即
\begin{equation*}
\int_{x^{0}>0}\mathrm{d}^{2}x\,b-\int_{x^{0}<0}\mathrm{d}^{2}x\,b=\frac{2\pi}{q}
\end{equation*}
这是与电荷量子 $q$ 相容的最小允许瞬子。在存在有限截断时，路径积分不仅应当包含规范场的光滑涨落，还应当包含瞬子涨落。

把瞬子效应包括进来之后，U(1) 规范理论就不再是自由理论。现在的问题是如何理解含瞬子的紧致 U(1) 理论的低能性质。一种方法是利用 XY 模型与 U(1) 规范理论之间的对偶性。

对虚时间而言，$\theta$ 与 $a_{\mu}$ 通过下式相联系：
\begin{equation*}
\frac{q}{2\pi}b=\frac{1}{\chi}\dot{\theta}, \qquad \frac{q}{2\pi}\epsilon_{ij}e_{j}=-\frac{1}{\chi}\partial_{i}\theta \tag{6.3.11}
\end{equation*}
此式也可写为
\begin{equation*}
\frac{q}{2\pi}\epsilon_{\mu\nu\lambda}\partial_{\nu}a_{\lambda}=\frac{1}{\chi}\partial_{\mu}\theta \tag{6.3.12}
\end{equation*}
把式 (6.3.12) 代入式 (6.3.10) 后，我们得到
\begin{equation*}
\mathcal{L}=\frac{\chi}{2}(\partial_{\mu}\theta)^{2}
\end{equation*}
其中 $\chi=\left(\frac{qg}{2\pi}\right)^{2}$。我们注意到，瞬子产生 $2\pi/q$ 的通量。如前所见，这样一份通量对应于 XY 模型中的一个粒子。因此，瞬子产生或湮灭一个粒子。在 XY 模型中，产生或湮灭一个粒子的算符正是 $\mathrm{e}^{\pm\mathrm{i}\theta}$。我们发现，位于 $x^{\mu}$ 处的瞬子对应于 $\mathrm{e}^{\mathrm{i}\theta(x^{\mu})}$。由于瞬子与反瞬子在路径积分中以相同的权重出现，我们可以通过在 XY 模型中加入一项 $\mathrm{e}^{\mathrm{i}\theta}+\mathrm{e}^{-\mathrm{i}\theta}=2\cos(\theta)$ 来把瞬子效应包括进来，如下所示：
\begin{equation*}
\mathcal{L}=\frac{\chi}{2}(\partial_{\mu}\theta)^{2}-K\cos(\theta)
\end{equation*}

当 $\chi$ 很大时，$\theta$ 在 $\theta=0$ 附近的涨落很小。我们可以把 $\mathcal{L}$ 近似为
\begin{equation*}
\mathcal{L}=\frac{\chi}{2}(\partial_{\mu}\theta)^{2}+\frac{1}{2}K\theta^{2}
\end{equation*}

把瞬子效应包括进来之后，$\partial_{\mu}\theta$ 之间（或 $(b,e_{i})$ 之间）的关联变成短程的。我们看到，瞬子效应在 1+2 维中为 U(1) 规范场打开了能隙。这里我们想强调，$K\cos(\theta)$ 项是一个相关微扰，无论 $K$ 多么小，U(1) 规范玻色子都会获得能隙。

如果规范玻色子是无能隙的，那么在 1+2 维中带电粒子之间将通过对数势相互作用。电荷为 $q$ 的粒子与电荷为 $-q$ 的粒子之间的势由下式给出：
\begin{equation*}
\pi\frac{g^{2}q^{2}}{(2\pi)^{2}}\ln r=\pi\chi\ln r
\end{equation*}

把瞬子效应包括进来之后，规范场获得了能隙。但瞬子效应会如何影响带电粒子之间的长程相互作用呢？我们注意到，电荷为 $q$ 的粒子在 XY 模型中由一个涡旋描述。在瞬子效应存在时，即在 $K\cos(\theta)$ 项存在时，涡旋与反涡旋之间的势随两个涡旋的间距线性增长（见图 6.6）。因此，瞬子效应把对数势变成了电荷之间的线性势。

我们已经证明，在 1+2 维中 XY 模型可以用紧致 U(1) 规范理论来描述；我们也已经证明，U(1) 规范理论中的瞬子效应会打开能隙。这两个结果似乎相互矛盾：具有无能隙 Nambu--Goldstone 模的 XY 模型，怎么会等价于有能隙的 U(1) 规范理论呢？然而我们意识到，瞬子会产生通量，而通量对应于粒子数，所以瞬子的存在意味着相应 XY 模型中的粒子数不守恒。于是一切都自洽了。在瞬子存在时，粒子数不守恒。紧致 U(1) 规范理论描述的是一个不具有 U(1) 对称性的 XY 模型。这解释了为什么瞬子效应对应于 XY 模型中的势项 $K\cos(\theta)$。反过来，如果 XY 模型具有 U(1) 对称性并保持粒子数守恒，那么相应的 U(1) 规范理论中就不允许存在瞬子。

%FIG{6.6}: {XY 模型中相距 $l$ 的涡旋--反涡旋对：$\theta$ 场在连接两个涡旋的连线的两侧必须跳变 $2\pi$。在 $-K\cos(\theta)$ 项存在时，通常的涡旋拟设 $\theta(x,y)=\arctan(x/y)$ 付出巨大的能量，其大小按 $l^{2}$ 增长。为使能量最小，非零的 $\theta$ 被限制在阴影区内。（虚线是 $\theta$ 场的等值线。）这导致两个涡旋之间的线性禁闭势。}

\subsection*{问题 6.3.3}

\textbf{含瞬子的 1+2 维 U(1) 规范理论与含 $\cos(\theta)$ 项的 XY 模型之间的对偶性}

\begin{enumerate}
\item 计算在 $x^{\mu}$ 处有一个瞬子、在 $y^{\mu}$ 处有一个反瞬子的 U(1) 规范理论的配分函数 $Z(x^{\mu},y^{\mu})$，精确到一个整体常数因子。
\item 利用上述结果，写出包含多瞬子气体贡献的 U(1) 规范理论配分函数的表达式。
\item 计算在 $x^{\mu}$ 处插入 $\mathrm{e}^{\mathrm{i}\theta}$、在 $y^{\mu}$ 处插入 $\mathrm{e}^{-\mathrm{i}\theta}$ 的 XY 模型的下列配分函数 $Z(x^{\mu},y^{\mu})$，精确到一个整体常数因子：
\begin{equation*}
Z(x^{\mu},y^{\mu})=\int\mathcal{D}\theta\,\mathrm{e}^{\mathrm{i}\theta(x^{\mu})}\mathrm{e}^{-\mathrm{i}\theta(y^{\mu})}\mathrm{e}^{-\int\mathrm{d}^{3}x^{\mu}\,\mathcal{L}_{XY}(\theta)}
\end{equation*}
\item 证明：含瞬子的 U(1) 配分函数与加了 $K\cos(\theta)$ 项的 XY 模型配分函数相一致（按 $K$ 的幂次展开之后）。
\end{enumerate}

\section{格点上的量子 U(1) 规范理论}

\subsection{格点 U(1) 规范理论的拉格朗日量}

\begin{keypoints}
\item 格点 U(1) 规范理论由链上的变量 $a_{\pmb{i}\pmb{j}}$ 与格点上的变量 $a_{0}(\pmb{i})$ 描述。
\item 紧致与非紧致的格点 U(1) 规范理论。
\item 电场 $e_{\pmb{i}\pmb{j}}$ 的通量。
\end{keypoints}

U(1) 规范理论也可以定义在格点上。U(1) 格点规范理论使我们能够在任意维度中研究强耦合极限与禁闭。为简单起见，这里我们考虑二维正方格子上的 U(1) 规范理论。标势 $a_{0}(\pmb{x})$ 定义在每个格点 $\pmb{i}$ 上，记作 $a_{0}(\pmb{i})$；矢势 $\pmb{a}(\pmb{x})$ 定义在每条链上，记作 $a_{\pmb{i}\pmb{j}}$。我们注意到，$a_{\pmb{i}\pmb{j}}$ 与 $a_{\pmb{j}\pmb{i}}$ 并不独立，它们通过 $a_{\pmb{i}\pmb{j}}=-a_{\pmb{j}\pmb{i}}$ 相联系。连续场与格点场通过下式相联系：
\begin{equation*}
a_{0}(\pmb{j})=a_{0}(\pmb{x})|_{\pmb{x}=l\pmb{j}}, \qquad a_{\pmb{j}_{1}\pmb{j}_{2}}=\int_{l\pmb{j}_{1}}^{l\pmb{j}_{2}}\mathrm{d}x^{i}a_{i}(\pmb{x}) \tag{6.4.1}
\end{equation*}
其中 $l$ 是晶格常数。描述 U(1) 格点规范理论的拉格朗日量是 $a_{0}(\pmb{i})$ 与 $a_{\pmb{i}\pmb{j}}$ 的下列函数：
\begin{equation*}
L=\frac{1}{4J}\sum_{\pmb{i},\pmb{\mu}=\pmb{x},\pmb{y}}\left[\dot{a}_{\pmb{i},\pmb{i}+\pmb{\mu}}+a_{0}(\pmb{i})-a_{0}(\pmb{i}+\pmb{\mu})\right]^{2}-\frac{g}{2}\sum_{\pmb{p}}\Phi_{\pmb{p}}^{2} \tag{6.4.2}
\end{equation*}
其中 $\Phi_{\pmb{p}}$ 是穿过标记为 $\pmb{p}$ 的正方形的 U(1) 通量：
\begin{equation*}
\Phi_{\pmb{p}}=a_{\pmb{i},\pmb{i}+\pmb{x}}+a_{\pmb{i}+\pmb{x},\pmb{i}+\pmb{x}+\pmb{y}}+a_{\pmb{i}+\pmb{x}+\pmb{y},\pmb{i}+\pmb{y}}+a_{\pmb{i}+\pmb{y},\pmb{i}}. \tag{6.4.3}
\end{equation*}
可以直接验证，上述拉格朗日量在下列格点 U(1) 规范变换
\begin{equation*}
a_{\pmb{i}\pmb{j}}(t)\rightarrow a_{\pmb{i}\pmb{j}}(t)+\phi_{\pmb{j}}(t)-\phi_{\pmb{i}}(t), \qquad a_{0}(\pmb{i},t)\rightarrow a_{0}(\pmb{i},t)+\dot{\phi}_{\pmb{i}}(t) \tag{6.4.4}
\end{equation*}
之下不变。

式 (6.4.2) 定义的是一个非紧致的格点 U(1) 规范理论。紧致的格点 U(1) 规范理论则通过把 $u_{\pmb{i}\pmb{j}}=\mathrm{e}^{\mathrm{i}a_{\pmb{i}\pmb{j}}}$ 取作链变量来定义。在紧致格点 U(1) 规范理论中，$a_{\pmb{i}\pmb{j}}$ 与 $a'_{\pmb{i}\pmb{j}}=a_{\pmb{i}\pmb{j}}+2\pi$ 被视为等价。紧致格点 U(1) 规范理论的拉格朗日量具有如下形式：
\begin{equation*}
L=\frac{1}{4J}\sum_{\pmb{i},\pmb{\mu}=\pmb{x},\pmb{y}}\left[\dot{a}_{\pmb{i},\pmb{i}+\pmb{\mu}}+a_{0}(\pmb{i})-a_{0}(\pmb{i}+\pmb{\mu})\right]^{2}+g\sum_{\pmb{p}}\cos(\Phi_{\pmb{p}}) \tag{6.4.5}
\end{equation*}
它由式 (6.4.2) 修改而来，以便与周期性条件 $a_{\pmb{i}\pmb{j}}\sim a_{\pmb{i}\pmb{j}}+2\pi$ 相容。在本节余下的部分中，我们只考虑紧致 U(1) 格点规范理论。

经典运动方程由 $\delta L/\delta a_{\pmb{i}\pmb{j}}=0$ 与 $\delta L/\delta a_{0}(\pmb{i})=0$ 得到。我们发现
\begin{equation*}
\frac{1}{2J}\frac{\mathrm{d}e_{\pmb{i},\pmb{i}+\pmb{\mu}}}{\mathrm{d}t}+g\left[\sin(\Phi_{\pmb{p}_{1}})-\sin(\Phi_{\pmb{p}_{2}})\right]=0, \tag{6.4.6}
\end{equation*}
\begin{equation*}
\sum_{\pmb{\mu}=\pm\pmb{x},\pm\pmb{y}}e_{\pmb{i},\pmb{i}+\pmb{\mu}}=0, \qquad e_{\pmb{i}\pmb{j}}=\dot{a}_{\pmb{i}\pmb{j}}+a_{0}(\pmb{i})-a_{0}(\pmb{j}) \tag{6.4.7}
\end{equation*}
其中 $\pmb{p}_{1}$ 与 $\pmb{p}_{2}$ 标记链 $(\pmb{i},\pmb{i}+\pmb{\mu})$ 两侧的两个正方形。

把 $e_{\pmb{i}\pmb{j}}$ 与连续模型中的 $e$ 作比较，我们看到 $e_{\pmb{i}\pmb{j}}$ 可以解释为流过链 $(\pmb{i},\pmb{j})$ 的电场通量。式 (6.4.7) 表明电通量是守恒的，它对应于连续模型中的高斯定律 $\pmb{\partial}_{\pmb{x}}\cdot\pmb{e}=0$。此外，$\Phi_{\pmb{p}}$ 作为穿过正方形 $S_{\pmb{p}}$ 的通量，对应于连续模型中的 $\Phi_{\pmb{p}}=\int_{S_{\pmb{p}}}\mathrm{d}^{2}\pmb{x}\,b$。

要获得式 (6.4.2) 或式 (6.4.5) 所描述的格点 U(1) 规范理论的动力学性质，最简单的办法是取连续极限。假设 $a_{\mu}$ 很小并且是 $t$ 与 $\pmb{x}$ 的光滑函数，我们把连续变量与格点变量之间的关系式 (6.4.1) 代入格点拉格朗日量，便得到下列标准的 U(1) 拉格朗日量：
\begin{equation*}
\mathcal{L}=\frac{1}{8\pi\alpha_{2D}}\left(\frac{1}{c}e^{2}-cb^{2}\right) \tag{6.4.8}
\end{equation*}
从所得的麦克斯韦方程我们发现，格点 U(1) 规范理论含有一个速度为 $c$ 的无能隙模式。耦合常数 $\alpha_{2D}$ 在 1+2 维中的量纲为 1/长度。

在 3.7.1 节中，我们曾利用规范不变性在连续空间中构造了带电玻色子与 U(1) 规范场相耦合的理论。同样的构造也适用于格点模型。我们发现，与带电玻色子耦合的格点 U(1) 规范理论由下列拉格朗日量描述：
\begin{align*}
L={}&+\sum_{\langle\pmb{i}\pmb{j}\rangle}t_{\pmb{i}\pmb{j}}\left(\varphi_{\pmb{i}}^{\dagger}\varphi_{\pmb{j}}\mathrm{e}^{-\mathrm{i}a_{\pmb{j}\pmb{i}}}+h.c\right)+\sum_{\pmb{i}}\varphi_{\pmb{i}}^{\dagger}\,\mathrm{i}\left[\partial_{0}+\mathrm{i}a_{0}(\pmb{i})\right]\varphi_{\pmb{i}}\\
&+\frac{1}{4J}\sum_{\pmb{i},\pmb{\mu}=\pmb{x},\pmb{y}}\left[\dot{a}_{\pmb{i},\pmb{i}+\pmb{\mu}}+a_{0}(\pmb{i})-a_{0}(\pmb{i}+\pmb{\mu})\right]^{2}+g\sum_{\pmb{p}}\cos(\Phi_{\pmb{p}}) \tag{6.4.9}
\end{align*}
该拉格朗日量在下列格点规范变换下不变：
\begin{equation*}
c_{\pmb{i}}\rightarrow\mathrm{e}^{\mathrm{i}\phi(\pmb{i},t)}c_{\pmb{i}}, \quad a_{0}(\pmb{i})\rightarrow a_{0}(\pmb{i})-\partial_{0}\phi(\pmb{i},t), \quad a_{\pmb{i}\pmb{j}}\rightarrow a_{\pmb{i}\pmb{j}}+\phi(\pmb{i},t)-\phi(\pmb{j},t).
\end{equation*}

\subsection*{问题 6.4.1}

\begin{enumerate}
\item[(a)] 把式 (6.4.1) 推广到三维立方格子。
\item[(b)] 利用式 (6.4.1) 证明，三维格点规范理论的拉格朗日量在连续极限下变为 $\frac{1}{8\pi\alpha}(c^{-1}\pmb{e}^{2}-c\pmb{b}^{2})$，其中 $e_{i}=\partial_{0}a_{i}-\partial_{i}a_{0}$，$b_{i}=\epsilon_{ijk}\partial_{j}a_{k}$。试用 $J$、$g$ 与晶格常数 $l$ 求出 $\alpha$ 和 $c$ 的值。这里无量纲常数 $\alpha$ 就是精细结构常数，$c$ 是光速。
\end{enumerate}

%FIG{6.7}: {一个正方形上的 U(1) 规范理论。}

\subsection{格点 U(1) 规范理论的哈密顿量}

\begin{keypoints}
\item 格点 U(1) 规范理论的哈密顿量可以通过固定规范或相空间路径积分得到。
\item 穿过一条链的电通量 $e_{\pmb{i}\pmb{j}}$ 是同一条链上的矢势 $a_{\pmb{i}\pmb{j}}$ 的正则“动量”。
\end{keypoints}

拉格朗日量 (6.4.5) 描述的是经典的 U(1) 格点规范理论。本节我们来求它的量子描述。为简单起见，我们考虑单个正方形上的格点规范理论（见图 6.7）。该量子理论由路径积分
\begin{equation*}
Z=\int\mathcal{D}a\,\mathcal{D}a_{0}\,\mathrm{e}^{\mathrm{i}\int\mathrm{d}t\left(\frac{1}{4J}\sum_{i}(\dot{a}_{i,i+1}+a_{0}(i)-a_{0}(i+1))^{2}+g\cos\Phi\right)}, \tag{6.4.10}
\end{equation*}
描述，其中 $i=1,2,3,4$。这里 $i=1$ 与 $i=5$ 被视为同一个点。作为一个规范理论，$(a_{ij}(t),a_{0}(i,t))$ 是路径的一种多对一标记。我们可以通过“固定规范”获得一对一的标记。我们注意到，$\sum_{j=i\pm1}a_{ij}$ 在规范变换下按
$\sum_{j=i\pm1}a_{ij}\rightarrow\sum_{j=i\pm1}\tilde{a}_{ij}=\sum_{j=i\pm1}(a_{ij}+\phi_{i}-\phi_{j})$
变换。通过调节 $\phi_{i}$，我们总能使 $\sum_{j=i\pm1}\tilde{a}_{ij}=0$。因此，对任意路径 $(a_{ij}(t),a_{0}(i,t))$，我们总能构造一个规范变换，使 $\sum_{j=i\pm1}a_{ij}=0$。于是，我们可以通过选择规范固定条件
\begin{equation*}
\sum_{j=i\pm1}a_{ij}=0
\end{equation*}
来固定规范。这样的规范称为库仑规范，对于连续理论它具有 $\pmb{\partial}\cdot\pmb{a}=0$ 的形式。在库仑规范下，我们的路径积分变为
\begin{equation*}
Z=\int\mathcal{D}a\,\mathcal{D}a_{0}\,\mathrm{e}^{-\mathrm{i}\int\mathrm{d}t\left(\frac{1}{4J}\sum_{i}(\dot{a}_{i,i+1}+a_{0}(i)-a_{0}(i+1))^{2}+g\cos\Phi\right)}\prod_{i,t}\delta\left(\sum_{j=i\pm1}a_{ij}\right)
\end{equation*}
我们注意到，$a_{0}(i)$ 与 $a_{ij}$ 之间的耦合具有 $a_{0}(i)\sum_{j=i\pm1}\dot{a}_{ij}$ 的形式。因此，对满足约束 $\sum_{j=i\pm1}a_{ij}=0$ 的 $a_{ij}$ 而言，$a_{0}(i)$ 与 $a_{ij}$ 解耦。由于 $a_{0}(i)$ 没有动力学（即不含 $\dot{a}_{0}(i)$ 项），我们可以积掉 $a_{0}(i)$，这相当于简单地把 $a_{0}$ 去掉。所得的路径积分变为
\begin{equation*}
Z=\int\mathcal{D}a\,\mathrm{e}^{-\mathrm{i}\int\mathrm{d}t\left(\frac{1}{4J}\sum_{i}\dot{a}_{i,i-1}^{2}+g\cos\Phi\right)}\prod_{i,t}\delta\left(\sum_{j=i\pm1}a_{ij}\right).
\end{equation*}
一般地，库仑规范下的路径积分可以由两个简单步骤得到：第一步，插入规范固定条件 $\prod_{i,t}\allowbreak\delta(\sum_{j=i\pm1}\allowbreak a_{ij})$；第二步，去掉 $a_{0}(i)$ 场。

对我们的问题，约束 $\prod_{i,t}\delta(\sum_{j=i\pm1}a_{ij})$ 要求 $a_{12}=a_{23}=a_{34}=a_{41}\equiv\Phi/4$。这里 $\Phi$ 描述穿过正方形的 U(1) 通量。路径积分取如下简单形式：
\begin{equation*}
Z=\int\mathcal{D}\theta\,\mathrm{e}^{-\mathrm{i}\int\mathrm{d}t\left(\frac{1}{16J}\dot{\Phi}^{2}+g\cos\Phi\right)} \tag{6.4.11}
\end{equation*}
我们注意到，位形 $(a_{12},a_{23},a_{34},a_{41})=(\pi/2,\pi/2,\pi/2,\pi/2)$ 与 $(a_{12},a_{23},a_{34},a_{41})=(2\pi,0,0,0)$ 规范等价（即存在一个把 $(\pi/2,\pi/2,\pi/2,\pi/2)$ 变到 $(2\pi,0,0,0)$ 的规范变换）。此外，由于 $a_{i,i+1}$ 生活在圆上，$a_{12}=2\pi$ 与 $a_{12}=0$ 等价。于是，$\Phi=2\pi$ 与 $\Phi=0$ 对应同一个物理点。路径积分 (6.4.11) 描述的是单位圆上一个质量为 $(8J)^{-1}$ 的粒子，通量能量 $-g\cos\Phi$ 就是该粒子所感受的势。当 $g=0$ 时，能级为 $E_{n}=4Jn^{2}$。

求解式 (6.4.10) 还有另一种方法。我们先引入 $a_{ij}$ 的正则动量，即
\begin{equation*}
\frac{\partial L}{\partial\dot{a}_{ij}}=\frac{1}{2J}e_{ij}
\end{equation*}
并把路径积分 (6.4.10) 写成相空间路径积分（见问题 6.4.4）
\begin{equation*}
Z=\int\mathcal{D}a\,\mathcal{D}a_{0}\,\mathcal{D}e_{ij}\,\mathrm{e}^{\mathrm{i}\int\mathrm{d}t\left(\sum_{i}\frac{e_{i,i+1}}{2J}(\dot{a}_{i,i+1}+a_{0}(i)-a_{0}(i+1))-\frac{1}{4J}\sum_{i}e_{i,i+1}^{2}+g\cos\Phi\right)}. \tag{6.4.12}
\end{equation*}

%FIG{6.8}: {带一条对角链的正方形上的 U(1) 规范理论。}

然后我们积掉 $a_{0}$，得到\footnote{注意 $a_{0}(i)$ 与 $e_{ij}$ 之间的耦合具有 $\sum_{i}a_{0}(i)(e_{i,i+1}+e_{i,i-1})$ 的形式。路径积分 $\int\mathcal{D}a_{0}\,\mathrm{e}^{\mathrm{i}\int\mathrm{d}t\sum_{i}a_{0}(i)(e_{i,i+1}+e_{i,i-1})}$ 正比于 $\prod_{i,t}\delta(e_{i,i+1}(t)+e_{i,i-1}(t))$。}
\begin{equation*}
Z=\int\mathcal{D}a\,\mathcal{D}e_{ij}\,\mathrm{e}^{\mathrm{i}\int\mathrm{d}t\left(\sum_{i}\frac{e_{i,i+1}}{2J}\dot{a}_{i,i+1}-\frac{1}{4J}\sum_{i}e_{i,i+1}^{2}+g\cos\Phi\right)}\prod_{i,t}\delta\left(e_{i,i+1}+e_{i,i-1}\right) \tag{6.4.13}
\end{equation*}
上述系统的量子哈密顿量为
\begin{equation*}
H=\frac{1}{4J}\sum_{i}e_{i,i+1}^{2}-g\cos\Phi, \qquad \left[a_{i,i+1},(2J)^{-1}e_{i,i+1}\right]=\mathrm{i}\hbar.
\end{equation*}
物理希尔伯特空间由满足
\begin{equation*}
\left(e_{i,i+1}+e_{i,i-1}\right)|\text{phy}\rangle=0
\end{equation*}
的态构成。由于 $\left[e_{i,i+1}+e_{i,i-1},H\right]=0$，$H$ 在物理希尔伯特空间内部起作用。

为求系统的能级，我们注意到 $a_{ij}$ 以 $2\pi$ 为周期，而 $a_{ij}$ 的动量是 $e_{ij}/2J$。于是，$e_{ij}/2J$ 量子化为整数 $n_{ij}$。当 $g=0$ 时，系统的能级为 $J\sum_{i}n_{i,i+1}^{2}$，其中 $n_{ij}$ 满足约束 $\sum_{j=i\pm1}n_{ij}=0$。对我们的四格点系统，该约束要求 $n_{12}=n_{23}=n_{34}=n_{41}\equiv n$。能量本征态由 $n$ 标记，能量为 $E_{n}=4Jn^{2}$，与先前的计算一致。

\subsection*{问题 6.4.2}

\begin{enumerate}
\item[(a)] 把式 (6.4.10) 推广到带一条对角链的正方形（见图 6.8）。
\item[(b)] 假设 $g=0$，求基态与激发态的能量。
\end{enumerate}

\subsection*{问题 6.4.3}

\begin{enumerate}
\item[(a)] 证明：对于二维正方格子上的非紧致格点 U(1) 规范理论，固定规范后的路径积分具有如下形式：
\begin{equation*}
Z=\int\mathcal{D}a\,\mathrm{e}^{-\mathrm{i}\int\mathrm{d}t\left(\frac{1}{4J}\sum_{i,\mu=x,y}\dot{a}_{i,i+\mu}^{2}-\frac{1}{2g}\sum_{p}\Phi_{p}^{2}\right)}\prod_{i,t}\delta\left(\sum_{\mu=\pm x,\pm y}a_{i,i+\mu}\right)
\end{equation*}
\item[(b)] 证明：约束 $\sum_{\mu=\pm x,\pm y}a_{i,i+\mu}=0$ 可以通过引入一个定义在正方形上的实场 $\phi_{p}$ 来求解，并把链上的规范势写成 $a_{i,i+\mu}=\phi_{p_{1}}-\phi_{p_{2}}$，其中 $p_{1}$ 与 $p_{2}$ 是链 $(i,i+\mu)$ 两侧的两个正方形。试用 $\phi_{p}$ 场写出路径积分。
\item[(c)] 求 $\phi_{p}$ 涨落的色散关系。
\end{enumerate}

\subsection*{问题 6.4.4}

证明式 (6.4.12)。（提示：相空间路径积分由 $\int\mathcal{D}p\,\mathcal{D}q\,\mathrm{e}^{\mathrm{i}\int\mathrm{d}t[p\dot{q}-H(p,q)]}$ 给出，其中 $H(p,q)=p\dot{q}-L$ 是哈密顿量。）

\subsection{格点 U(1) 规范理论的库仑相与禁闭相}

\begin{keypoints}
\item 在 1+3 维中，当 $g/J\gg1$ 时，紧致格点 U(1) 规范理论处于库仑相；当 $g/J\ll1$ 时，处于禁闭相。
\end{keypoints}

我们知道，在 1+2 维中，由于瞬子效应，紧致 U(1) 规范理论总是处于禁闭相。本节我们将研究立方格子上的紧致 U(1) 规范理论。我们将论证，1+3 维的 U(1) 规范理论既有禁闭相，也有库仑相。

U(1) 规范理论的拉格朗日量为
\begin{equation*}
L=\frac{1}{4J}\sum_{\pmb{i},\pmb{\mu}=\pmb{x},\pmb{y},\pmb{z}}\left[\dot{a}_{\pmb{i},\pmb{i}+\pmb{\mu}}+a_{0}(\pmb{i})-a_{0}(\pmb{i}+\pmb{\mu})\right]^{2}+g\sum_{\pmb{p}}\cos(\Phi_{\pmb{p}})
\end{equation*}
我们可以把上述拉格朗日量的作用量写成下列无量纲形式：
\begin{equation*}
S=\sqrt{\frac{g}{J}}\int\mathrm{d}\tilde{t}\left(\frac{1}{4}\sum_{\pmb{i},\pmb{\mu}=\pmb{x},\pmb{y}}\left[\partial_{\tilde{t}}a_{\pmb{i},\pmb{i}+\pmb{\mu}}+\tilde{a}_{0}(\pmb{i})-\tilde{a}_{0}(\pmb{i}+\pmb{\mu})\right]^{2}+\sum_{\pmb{p}}\cos(\Phi_{\pmb{p}})\right)
\end{equation*}
（译注：原文此处求和指标为 $\pmb{\mu}=\pmb{x},\pmb{y}$，按上下文疑为 $\pmb{x},\pmb{y},\pmb{z}$ 之误。）其中 $\tilde{t}=\sqrt{gJ}\,t$ 是无量纲时间，$\tilde{a}_{0}=a_{0}/\sqrt{gJ}$ 是无量纲势。我们看到，当 $g/J\gg1$ 时，$\Phi_{\pmb{p}}$ 的涨落很弱，我们可以假定 $a_{\pmb{i}\pmb{j}}\sim0$。在这个极限下，$L$ 变成如下非紧致的 U(1) 规范理论：
\begin{equation*}
L=\frac{1}{4J}\sum_{\pmb{i},\pmb{\mu}=\pmb{x},\pmb{y},\pmb{z}}\left[\dot{a}_{\pmb{i},\pmb{i}+\pmb{\mu}}+a_{0}(\pmb{i})-a_{0}(\pmb{i}+\pmb{\mu})\right]^{2}-\frac{g}{2}\sum_{\pmb{p}}\Phi_{\pmb{p}}^{2}
\end{equation*}
这个二次型理论可以精确求解。在连续极限下，上式变为 U(1) 规范理论的标准麦克斯韦拉格朗日量。我们发现，低能下存在两个以 $\omega=c_{a}|\pmb{k}|$ 线性色散的模式，它们就是无能隙的光子。

当 $g/J\ll1$ 时，我们可以去掉 $\cos(\Phi_{\pmb{p}})$ 项。所得的二次型理论同样可以精确求解，就像上一节中所做的那样。相空间路径积分的拉格朗日量具有如下形式：
\begin{equation*}
L=\sum_{\pmb{i},\pmb{\mu}=\pmb{x},\pmb{y},\pmb{z}}\frac{e_{\pmb{i},\pmb{i}+\pmb{\mu}}}{2J}\left[\dot{a}_{\pmb{i},\pmb{i}+\pmb{\mu}}+a_{0}(\pmb{i})-a_{0}(\pmb{i}+\pmb{\mu})\right]-\sum_{\pmb{i},\pmb{\mu}=\pmb{x},\pmb{y},\pmb{z}}\frac{e_{\pmb{i},\pmb{i}+\pmb{\mu}}^{2}}{4J}+g\sum_{\pmb{p}}\cos(\Phi_{\pmb{p}})
\end{equation*}
积掉 $a_{0}$ 之后，我们得到约束
\begin{equation*}
\sum_{\pmb{\mu}=\pm\pmb{x},\pm\pmb{y},\pm\pmb{z}}e_{\pmb{i},\pmb{i}+\pmb{\mu}}=0
\end{equation*}
哈密顿量为
\begin{equation*}
H=\sum_{\pmb{i},\pmb{\mu}=\pmb{x},\pmb{y},\pmb{z}}\frac{e_{\pmb{i},\pmb{i}+\pmb{\mu}}^{2}}{4J}-g\sum_{\pmb{p}}\cos(\Phi_{\pmb{p}})
\end{equation*}
由于 $e_{\pmb{i}\pmb{j}}/2J$ 量子化为整数 $n_{\pmb{i}\pmb{j}}$，当 $g=0$ 时，能级为 $E=J\sum_{\pmb{i},\pmb{\mu}=\pmb{x},\pmb{y},\pmb{z}}n_{\pmb{i},\pmb{i}+\pmb{\mu}}^{2}$，其中 $n_{\pmb{i},\pmb{i}+\pmb{\mu}}$ 满足 $\sum_{\pmb{\mu}=\pm\pmb{x},\pm\pmb{y},\pm\pmb{z}}n_{\pmb{i},\pmb{i}+\pmb{\mu}}=0$。我们看到，$g=0$ 极限下的所有激发都是有能隙的。这些有能隙的激发是由非零的 $n_{\pmb{i}\pmb{j}}$ 构成的圈，这些圈代表电通量线。当存在一对正、负电荷时，两个电荷由一条电通量线连接，这在两个电荷之间产生线性禁闭势。因此，这个有能隙的相就是 U(1) 规范理论的禁闭相。

\subsection*{问题 6.4.5}

为证明在禁闭相中电荷之间以线性势相互作用，考虑存在电荷时的格点 U(1) 规范理论：
\begin{equation*}
L=\sum_{\pmb{i},\pmb{\mu}=\pmb{x},\pmb{y},\pmb{z}}\frac{\left[\dot{a}_{\pmb{i},\pmb{i}+\pmb{\mu}}+a_{0}(\pmb{i})-a_{0}(\pmb{i}+\pmb{\mu})\right]^{2}}{4J}+g\sum_{\pmb{p}}\cos(\Phi_{\pmb{p}})-\sum a_{0}(\pmb{i})q_{\pmb{i}}
\end{equation*}
其中 $q_{\pmb{i}}$ 是格点 $\pmb{i}$ 上的 U(1) 电荷。
\begin{enumerate}
\item[(a)] 求相空间路径积分中的拉格朗日量。
\item[(b)] 证明：积掉 $a_{0}$ 导出如下约束：
\begin{equation*}
\sum_{\pmb{\mu}=\pm\pmb{x},\pm\pmb{y},\pm\pmb{z}}e_{\pmb{i},\pmb{i}+\pmb{\mu}}=2Jq_{\pmb{i}}
\end{equation*}
这就是格点上的高斯定律。
\item[(c)] 假设 $g=0$，求在 $\pmb{i}=0$ 处有电荷 $q=1$、在 $\pmb{i}=l\pmb{x}$ 处有电荷 $q=-1$ 的规范理论的基态能量。证明两个电荷之间的相互作用能随 $l$ 线性增长。
\end{enumerate}
